% !PRINTED 546
\weberpage{546}
If $\omega$ belongs to the exponent $a$, then the powers
\weq{9}{1, \omega, \omega^2, \ldots, \omega^{a-1}}
are all incongruent, and therefore, since they all satisfy the congruence (8) (by 2.), they form the totality of the roots of this congruence. Among the numbers (9), therefore, all other numbers $\omega$ belonging to the exponent $a$ must be sought. But $\omega^l$ will belong to the exponent $a$ only when $l$ is relatively prime to $a$, and it follows:

If there are any numbers $\omega$ at all that belong to the exponent $a$, then their number is as great as the number of relative primes to $a$ in the sequence of numbers $0, 1, 2, \ldots, a-1$. We denote this number, as we did earlier (Bd. I, \S.\,132), by $\varphi(a)$. But that to every divisor $a$ of $p^f-1$ there always belongs at least one number $\omega$, and consequently $\varphi(a)$ numbers, can be proved in just the same way as the corresponding theorem of rational number theory is proved in \S.\,136 of the first volume.

The numbers $\omega$ that belong to the exponent $p^f-1$, of which by the foregoing there are always $\varphi(p^f-1)$, are called \gesperrt{primitive roots} of $\pi$. If $\gamma$ is such a primitive root, then the powers
\[ 1, \gamma, \gamma^2, \ldots, \gamma^{p^f-2} \]
form a complete system of residues to the modulus $\pi$, with the exclusion of the number divisible by $\pi$.

By \gesperrt{Fermat}'s theorem for rational numbers, $t^{p-1}-1$ is divisible by $p$, and consequently also by $\pi$, for $t = 1, 2, \ldots, p-1$. The congruence
\weq{10}{t^{p-1} \equiv 1 \pmod{\pi}}
therefore has the roots
\[ 1, 2, \ldots, p-1, \]
and these are incongruent with one another, since by \S.\,140,3. a rational number is divisible by $\pi$ only when it is divisible by $p$. By Theorem 2., therefore, the congruence (10) has no roots other than these.

If we multiply the congruence (10) by $t$ as well, it follows that the congruence of degree $p$
\[ t^p - t \equiv 0 \pmod{\pi} \]
has the $p$ roots
\[ 0, 1, 2, \ldots, p-1, \]
%CONT>>
